\section{Modelo Conceitual}

\subsection{Diagrama Entidade-Relacionamento}

O modelo conceitual representa as entidades principais e seus relacionamentos:

\subsubsection{Entidades Principais}

\paragraph{Usuario}
\begin{itemize}
    \item Pessoa fisica que utiliza o sistema
    \item Pode estar associado a um grupo (casal/familia)
    \item Possui configuracoes individuais
\end{itemize}

\paragraph{Grupo}
\begin{itemize}
    \item Agrupamento de usuarios (casal, familia, republica)
    \item Define o modelo de organizacao financeira
    \item Compartilha dados conforme configuracao
\end{itemize}

\paragraph{Conta}
\begin{itemize}
    \item Representa contas bancarias, carteiras, investimentos
    \item Pode ser individual ou compartilhada
    \item Possui saldo atual e historico
\end{itemize}

\paragraph{Categoria}
\begin{itemize}
    \item Classificacao de receitas e despesas
    \item Pode ser personalizada por usuario/grupo
    \item Possui subcategorias hierarquicas
\end{itemize}

\paragraph{Transacao}
\begin{itemize}
    \item Movimentacao financeira (entrada ou saida)
    \item Associada a uma conta e categoria
    \item Possui informacoes de divisao (para modelo Splitwise)
\end{itemize}

\paragraph{Orcamento}
\begin{itemize}
    \item Planejamento financeiro por categoria/periodo
    \item Pode ser individual ou compartilhado
    \item Possui metas e limites de gasto
\end{itemize}

\paragraph{Meta}
\begin{itemize}
    \item Objetivos de poupanca ou investimento
    \item Possui valor alvo e prazo
    \item Acompanha progresso automaticamente
\end{itemize}

\subsection{Relacionamentos}

\subsubsection{Usuario - Grupo (N:M)}
\begin{itemize}
    \item Um usuario pode pertencer a multiplos grupos
    \item Um grupo pode ter multiplos usuarios
    \item Relacionamento possui papel (principal, secundario, dependente)
\end{itemize}

\subsubsection{Usuario - Conta (1:N)}
\begin{itemize}
    \item Um usuario pode ter multiplas contas
    \item Uma conta pertence a um usuario especifico
    \item Contas compartilhadas referenciam o grupo
\end{itemize}

\subsubsection{Conta - Transacao (1:N)}
\begin{itemize}
    \item Uma conta pode ter multiplas transacoes
    \item Uma transacao pertence a uma conta especifica
\end{itemize}

\subsubsection{Categoria - Transacao (1:N)}
\begin{itemize}
    \item Uma categoria pode ter multiplas transacoes
    \item Uma transacao possui uma categoria principal
\end{itemize}

\subsubsection{Usuario - Orcamento (1:N)}
\begin{itemize}
    \item Um usuario pode ter multiplos orcamentos
    \item Um orcamento pode ser individual ou de grupo
\end{itemize}

\subsection{Cardinalidades e Restricoes}

\begin{itemize}
    \item \textbf{Usuario}: Obrigatorio ter pelo menos um email valido
    \item \textbf{Grupo}: Minimo 2 usuarios, maximo configuravel
    \item \textbf{Conta}: Deve ter saldo inicial definido
    \item \textbf{Transacao}: Valor deve ser diferente de zero
    \item \textbf{Categoria}: Nome deve ser unico por usuario/grupo
    \item \textbf{Orcamento}: Periodo deve ser valido (inicio < fim)
    \item \textbf{Meta}: Valor alvo deve ser positivo
\end{itemize}

\subsection{Regras de Negocio}

\subsubsection{Modelos de Organizacao}
\begin{enumerate}
    \item \textbf{Individual}: Todas as entidades sao privadas do usuario
    \item \textbf{Renda Conjunta}: Contas e orcamentos sao compartilhados
    \item \textbf{Splitwise}: Transacoes individuais, divisao calculada
    \item \textbf{Fundo Comum}: Hibrido entre individual e compartilhado
\end{enumerate}

\subsubsection{Integridade Referencial}
\begin{itemize}
    \item Exclusao de usuario: Transferir dados para outro membro do grupo
    \item Exclusao de grupo: Converter contas compartilhadas em individuais
    \item Exclusao de categoria: Mover transacoes para categoria "Outros"
    \item Exclusao de conta: Transferir saldo para conta principal
\end{itemize}
